% Source printed pp. 236-242 / PDF pp. 244-250.
\section{所有曲线的叠}\label{src:6.4}
\begin{epigraph}
在包含一些糟糕对象的好范畴中工作, 胜过在只包含好对象的糟糕范畴中工作.
\epigraphauthor{John C. Baez}
\end{epigraph}
我们证明所有 $n$-点固有曲线的叠 $\mathcal M_{g,\allowbreak n}^{\mathrm{all}}$ 是代数叠 (定理~\bookref{6.4.6}), 稳定曲线的开子叠 $\overline{\mathcal M}_{g,\allowbreak n}\subseteq\mathcal M_{g,\allowbreak n}^{\mathrm{all}}$ 是 $\Spec\mathbb Z$ 上光滑, 维数为 $3g-3+n$ 的拟紧 Deligne--Mumford 叠 (定理~\bookref{6.4.16}).
\subsection{任意曲线族}\label{src:6.4.1-subsection}
回顾域 $k$ 上的曲线是一维有限型 $k$-概形, 射影曲线 $C$ 的亏格为 $g(C)=1-\chi(C,\allowbreak \mathcal O_C)$. 曲线可以不是纯维的, 不连通, 且具有任意奇点 (例如非约化). 例如, 一条亏格为 $g$ 的曲线与 $g$ 个 $\mathbb P^1$ 的不交并的亏格为 $0$. 虽然有参数化几何约化曲线的代数叠 (注记~\bookref{6.4.20}), 一旦允许非约化曲线, 就必须也允许非纯维曲线, 否则该叠将不是代数叠. 例如, 由有理正规曲线 $\mathbb P^1\hookrightarrow\mathbb P^3$ 的平坦退化产生的结点处带嵌入点的平面结点曲线~\cite[例 9.8.4]{Har77}, 可在平坦族中形变为平面结点三次曲线与 $\mathbb P^3$ 中一个点的不交并.

概形 $S$ 上的 $n$-点曲线族是代数空间的固有, 平坦有限表现态射 $C\to\allowbreak S$, 配备 $n$ 个截面 $\sigma_1,\allowbreak \ldots,\allowbreak \sigma_n:S\to\allowbreak C$, 使每个纤维 $C_s$ 是曲线. 若每个纤维 $C_s$ 的亏格都是 $g$, 则称族 $(C\to\allowbreak S,\allowbreak \sigma_i)$ 的亏格为 $g$. 注意标记点 $\sigma_i(s)$ 在纤维 $C_s$ 中可以重合, 也可以是奇点. 为使曲线的叠成为代数叠, 必须允许总空间 $C$ 为代数空间 (见注意~\bookref{6.3.21}).
\begin{proposition}\label{src:6.4.1}
若 $C\to\allowbreak S$ 是概形 $S$ 上的曲线族, 则存在平展覆盖 $S'\to\allowbreak S$, 使得 $C_{S'}\to\allowbreak S'$ 射影.
\end{proposition}
\begin{proof}
给出从局部到整体的论证. 先归结到 $S$ 在 $\Spec\mathbb Z$ 上有限型的情形. 由 Noether 逼近 (\bookref{B.3.2}), 可将 $S=\varprojlim_{\lambda\in\Lambda}S_\lambda$ 写成有限表现 $\mathbb Z$-概形的极限, 转移映射仿射. 由态射沿极限的下降 (\bookref{B.3.3}), 存在指标 $0\in\Lambda$ 和有限表现态射 $C_0\to\allowbreak S_0$, 使 $C\cong C_0\times_{S_0}S_\lambda$; 而且若对 $\lambda\geq0$ 置 $C_\lambda=C_0\times_{S_0}S$, 则 $C=\varprojlim_{\lambda\geq0}C_\lambda$. 由态射性质沿极限的下降 (\bookref{B.3.7}), 对 $\lambda\gg0$, $C_\lambda\to\allowbreak S_\lambda$ 是曲线族. 因为射影性在基变换下稳定, 可将 $S$ 换为 $S_\lambda$.

对点 $s\in S$, 定义 $S_n=\Spec\mathcal O_{S,\allowbreak s}/\mathfrak m_s^{n+1}$ 和 $\widehat S=\Spec\widehat{\mathcal O}_{S,\allowbreak s}$. 考虑 Cartesian 图
\[
\begin{tikzcd}[column sep=small]
C_s=C_0\arrow[r]\arrow[d]&C_1\arrow[r]\arrow[d]&\cdots\arrow[r]&\widehat C\arrow[r]\arrow[d]&C\arrow[d]\\
\Spec\kappa(s)=S_0\arrow[r]&S_1\arrow[r]\arrow[ur,phantom,"\square"]&\cdots\arrow[r]&\widehat S\arrow[r]\arrow[ur,phantom,"\square"]&S.
\end{tikzcd}
\]
由于分离的一维代数空间是概形 (定理~\bookref{5.5.32}), 固有一维概形是射影的, $C_0$ 上存在丰沛线丛 $L_0$. 由命题~\bookref{C.2.11}, 将 $C_n$ 上的线丛 $L_n$ 形变到 $C_{n+1}$ 上的 $L_{n+1}$ 的障碍位于 $\mathrm H^2(C_0,\allowbreak \mathcal O_{C_0}\otimes_{\kappa(s)}\mathfrak m^n)$ 中. 因为 $\dim C_0=1$, 障碍消失, 因而存在 $C_n$ 上相容的线丛序列 $L_n$. 由 Grothendieck 存在定理 (\bookref{C.5.3}), 存在 $\widehat C$ 上延拓 $L_n$ 的线丛 $\widehat L$.

对函子 $\Sch/S\to\Sets$, $(T\to\allowbreak S)\mapsto\operatorname{Pic}(C_T)$ 应用 Artin 逼近 (定理~\bookref{B.5.18}), 得到 $s$ 的平展邻域 $(S',\allowbreak s')\to(S,\allowbreak s)$ 及 $C_{S'}$ 上延拓 $L_0$ 的线丛 $L'$. 因为 $L_0$ 丰沛, 且丰沛性在族中是开条件 (命题~\bookref{B.2.10}), 将 $S'$ 换为 $s'$ 的开邻域后, 可使 $L'$ 在 $S'$ 上相对丰沛.
\end{proof}
\begin{exercise}\label{src:6.4.2}
通过明确延拓丰沛线丛给出另一论证:
\begin{enumerate}[(i)]
\item 用极限方法归结到 $S$ 是严格 Hensel 局部环 $R$ 的谱, 闭点为 $s$, 并将陈述归结到 $\operatorname{Pic}(C)\to\operatorname{Pic}(C_s)$ 的满射性.
\item 用 Gabber 引理 (\bookref{5.6.1}) 归结到 $C$ 是概形的情形.
\item 归结到 $C$ 约化的情形, 并对每个闭点 $x_s\in C_s$ 明确延拓线丛 $\mathcal O_{C_s}(-x_s)$, 说明此时 $\operatorname{Pic}(C)\to\operatorname{Pic}(C_s)$ 满射.
\end{enumerate}
另见~\cite[IV.4.1]{SGA4half}, \cite[引理 1.2]{Hal13} 和~\cite[推论 13.2.5]{Ols16}.
\end{exercise}
\begin{remark}\label{src:6.4.3}
Raynaud 给出了仿射曲线上一个光滑亏格 $g=1$ 曲线族的例子, 它 Zariski 局部射影, 但不是射影的~\cite[XIII.3.1]{Ray70}. 注意~\bookref{6.3.21} 中的例子甚至不是 Zariski 局部射影的.
\end{remark}
\subsection{所有曲线的叠的代数性}\label{src:6.4.2-subsection}
\begin{definition}\label{src:6.4.4}
设 $\mathcal M_{g,\allowbreak n}^{\mathrm{all}}$ 是 $\Sch_{\acute et}$ 上的预叠, 概形 $S$ 上的对象是具有 $n$ 个截面 $\sigma_1,\allowbreak \ldots,\allowbreak \sigma_n:S\to\allowbreak C$ 的亏格 $g$ 曲线族 $C\to\allowbreak S$. 态射 $(C'\to\allowbreak S',\allowbreak \sigma'_1,\allowbreak \ldots,\allowbreak \sigma'_n)\to(C\to\allowbreak S,\allowbreak \sigma_1,\allowbreak \ldots,\allowbreak \sigma_n)$ 是 Cartesian 图
\[
\begin{tikzcd}
C'\arrow[r,"g"]\arrow[d,shift right=1ex]&C\arrow[d,shift right=1ex]\\
S'\arrow[r,"f"']\arrow[u,shift right=1ex,"\sigma'_i"]&S\arrow[u,shift right=1ex,"\sigma_i"']
\end{tikzcd}
\]
的数据, 满足 $g\circ\sigma'_i=\sigma_i\circ f$.
\end{definition}
作为证明 $\mathcal M_{g,\allowbreak n}^{\mathrm{all}}$ 代数性的准备, 先说明对角线可表.
\begin{lemma}\label{src:6.4.5}
对角线 $\mathcal M_{g,\allowbreak n}^{\mathrm{all}}\to\mathcal M_{g,\allowbreak n}^{\mathrm{all}}\times\mathcal M_{g,\allowbreak n}^{\mathrm{all}}$ 可表.
\end{lemma}
\begin{proof}
为简单起见, 处理 $n=0$ 的情形. 设 $S$ 是概形, 态射 $S\to\mathcal M_g^{\mathrm{all}}\times\mathcal M_g^{\mathrm{all}}$ 对应于曲线族 $C_1\to\allowbreak S$ 和 $C_2\to\allowbreak S$. 考虑 Cartesian 图
\[
\begin{tikzcd}
\underline{\operatorname{Isom}}_S(C_1,C_2)\arrow[r]\arrow[d]&\mathcal M_g^{\mathrm{all}}\arrow[d]\\
S\arrow[r]&\mathcal M_g^{\mathrm{all}}\times\mathcal M_g^{\mathrm{all}},
\end{tikzcd}
\]
须说明 $\underline{\operatorname{Isom}}_S(C_1,\allowbreak C_2)$ 是代数空间. 由命题~\bookref{6.4.1}, 存在平展覆盖 $S'\to\allowbreak S$, 使得 $C_{S'}\to\allowbreak S'$ 射影. 由于 $\underline{\operatorname{Isom}}_S(C_1,\allowbreak C_2)\times_S S'=\underline{\operatorname{Isom}}_{S'}(C_{1,\allowbreak S'},\allowbreak C_{2,\allowbreak S'})$, 态射 $\underline{\operatorname{Isom}}_{S'}(C_{1,\allowbreak S'},\allowbreak C_{2,\allowbreak S'})\to\underline{\operatorname{Isom}}_S(C_1,\allowbreak C_2)$ 满射, 平展且可表. 因而可假设 $C_1,\allowbreak C_2$ 在 $S$ 上射影.

考虑函子的包含
\[
\underline{\operatorname{Isom}}_S(C_1,C_2)\subseteq\underline{\operatorname{Mor}}_S(C_1,C_2)\subseteq\Hilb(C_1\times_S C_2/S),
\]
第二个包含将态射 $\alpha:C_1\to\allowbreak C_2$ 映到其图 $\Gamma_\alpha:C_1\hookrightarrow C_1\times_S C_2$ ($T$-值点上的定义类似). 由于参数化使 $(C_1)_S\to(C_2)_S$ 为同构的映射 $T\to\allowbreak S$ 的 $\operatorname{Mor}(-,\allowbreak S)$ 的子函子是开的 (习题~\bookref{1.3.46}), 第一个包含是可表开浸入. 对第二个包含, 子概形 $\{Z\subseteq C_1\times_S C_2\}\in\Hilb(C_1\times_S C_2/S)(S)$ 位于 $\operatorname{Mor}(C_1,\allowbreak C_2)(S)$ 的像中当且仅当复合 $Z\hookrightarrow C_1\times_S C_2\xrightarrow{p_1}C_1$ 是同构 ($T$-值点同理). 因此习题~\bookref{1.3.46} 也说明第二个包含是可表开浸入.
\end{proof}
\begin{theorem}\label{src:6.4.6}
$\mathcal M_{g,\allowbreak n}^{\mathrm{all}}$ 是 $\Spec\mathbb Z$ 上局部有限型的代数叠.
\end{theorem}
\begin{proof}
为说明 $\mathcal M_{g,\allowbreak n}^{\mathrm{all}}$ 是 $\Sch_{\acute et}$ 上的叠, 设 $\{S_i\to\allowbreak S\}$ 是概形的平展覆盖, 对每个 $i$, $(C_i\to\allowbreak S_i,\allowbreak \sigma_{i,\allowbreak 1},\allowbreak \ldots,\allowbreak \sigma_{i,\allowbreak n})$ 是 $n$-点曲线族, 且 $\alpha_{ij}:C_i|_{S_{ij}}\to\allowbreak C_j|_{S_{ij}}$ 是 $S_{ij}=S_i\times_S S_j$ 上与截面相容且满足上循环条件的同构. 由 $p_1,\allowbreak p_2\circ\alpha_{ij}$ 定义映射的平展等价关系 $\coprod_{i,\allowbreak j}C_{ij}\rightrightarrows\coprod_i C_i$ 的商是代数空间 $C$. 此外, 由态射的平展下降, 有截面 $\sigma_1,\allowbreak \ldots,\allowbreak \sigma_n:S\to\allowbreak C$, 满足 $\sigma_{k,\allowbreak i}=\sigma_k|_{S_i}$. 因而 $(C\to\allowbreak S,\allowbreak \sigma_k)$ 是在每个 $S_i$ 上限制为 $(C_i\to\allowbreak S_i,\allowbreak \sigma_{i,\allowbreak k})$ 的 $n$-点曲线族.

为说明代数性, 断言只需处理 $n=0$. 忘掉最后一个标记点定义的映射 $\mathcal M_{g,\allowbreak n+1}^{\mathrm{all}}\to\mathcal M_{g,\allowbreak n}^{\mathrm{all}}$ 良定义且可表: 若 $S\to\mathcal M_{g,\allowbreak n}^{\mathrm{all}}$ 对应于 $n$-点族 $(C\to\allowbreak S,\allowbreak \sigma_k)$, 则 $S\times_{\mathcal M_{g,\allowbreak n}^{\mathrm{all}}}\mathcal M_{g,\allowbreak n+1}^{\mathrm{all}}\cong C$. (事实上, 该映射等同于万有族; 见习题~\bookref{6.4.8}.) 因此若 $U\to\mathcal M_{g,\allowbreak n}^{\mathrm{all}}$ 是由概形给出的光滑呈示, 则 $U':=U\times_{\mathcal M_{g,\allowbreak n}^{\mathrm{all}}}\mathcal M_{g,\allowbreak n+1}^{\mathrm{all}}\to\mathcal M_{g,\allowbreak n+1}^{\mathrm{all}}$ 是来自代数空间 $U'$ 的满射, 平展且可表的映射. 选取概形给出的平展呈示 $V\to\allowbreak U'$, 复合 $V\to\allowbreak U'\to\mathcal M_{g,\allowbreak n+1}^{\mathrm{all}}$ 是光滑呈示.

为说明 $\mathcal M_g^{\mathrm{all}}$ 代数, 只需说明对域 $k$ 上的每条射影曲线 $X$, 存在来自 $\mathbb Z$ 上有限型概形 $U$ 的光滑可表态射 $U\to\mathcal M_g^{\mathrm{all}}$, 其像包含 $[X]\in|\mathcal M_g^{\mathrm{all}}|$. 取嵌入 $X\hookrightarrow\mathbb P^N$, 使 $\mathrm H^1(X,\allowbreak \mathcal O_X(1))=0$, 设 $P(t)$ 是其 Hilbert 多项式. 设 $H:=\Hilb^P(\mathbb P^N_{\mathbb Z}/\mathbb Z)$ 为 Hilbert 概形, 它由定理~\bookref{2.1.2} 在 $\mathbb Z$ 上射影. 考虑万有族
\[
\begin{tikzcd}
C\arrow[r,hook]\arrow[d]&\mathbb P^N_H\arrow[dl]\\
H,&
\end{tikzcd}
\]
存在点 $h\in H(k)$, 使得作为 $\mathbb P^N_k$ 的闭子概形有 $C_h=X$. 上同调与基变换 (\bookref{A.6.8}) 说明 $h$ 有开邻域 $H'\subseteq H$, 对所有 $s\in H'$, $\mathrm H^1(C_s,\allowbreak \mathcal O_{C_s}(1))=0$. 考虑忘掉嵌入定义的态射 $H'\to\mathcal M_g^{\mathrm{all}}$, $[C\hookrightarrow\mathbb P^n]\mapsto[C]$. 由 $\mathcal M_g^{\mathrm{all}}$ 的对角线可表 (引理~\bookref{6.4.5}), $H'\to\mathcal M_g^{\mathrm{all}}$ 可表 (推论~\bookref{4.2.3}).

断言 $H'\to\mathcal M_g^{\mathrm{all}}$ 光滑. 虽然尚未证明 $\mathcal M_g^{\mathrm{all}}$ 代数, 但已知该态射可表, 这足以应用无穷小提升判据 (\bookref{4.7.1}) 检验光滑性. 设 $A'\to\allowbreak A$ 是剩余域为 $k$ 的 Artin 局部环的满射, 满足 $k=\ker(A'\to\allowbreak A)$. 对每条嵌入曲线 $[C\subseteq\mathbb P^N]\in H'$ 及图
\begin{equation}\label{src:6.4.7}
\begin{tikzcd}[column sep=large]
\Spec k\arrow[r]\arrow[rr,bend left=35,"{[C\subseteq\mathbb P^N_k]}"]&\Spec A\arrow[r,"{[\mathcal C\subseteq\mathbb P^N_A]}"]\arrow[d]&H'\arrow[d]\\
&\Spec A'\arrow[r,"{[\mathcal C']}"']\arrow[ur,dashed,"{[\mathcal C'\subseteq\mathbb P^N_{A'}]}"']&\mathcal M_g^{\mathrm{all}},
\end{tikzcd}
\end{equation}
须说明存在延拓此图的虚线箭头. 这等价于以下图中虚线箭头的存在:
\[
\begin{tikzcd}[column sep=small,row sep=small]
&\mathbb P^N_k&&\mathbb P^N_A&&\mathbb P^N_{A'}\arrow[ddl]\\
C\arrow[ur,hook]\arrow[rr]\arrow[d]&&\mathcal C\arrow[ur,hook]\arrow[rr]\arrow[d]&&\mathcal C'\arrow[ur,dashed]\arrow[d]&\\
\Spec k\arrow[rr]&&\Spec A\arrow[rr]&&\Spec A'.&
\end{tikzcd}
\]
闭浸入 $\mathcal C\hookrightarrow\mathbb P^N_A$ 由 $\mathcal C$ 上的甚丰沛线丛 $L$ 和截面 $s_0,\allowbreak \ldots,\allowbreak s_N\in\Gamma(\mathcal C,\allowbreak L)$ 定义. 将 $L$ 形变到 $\mathcal C'$ 的障碍位于 $\mathrm H^2(C,\allowbreak \mathcal O_C)=0$ 中 (命题~\bookref{C.2.11}), 故可将 $L$ 延拓为 $\mathcal C'$ 上的线丛 $L'$. 现在说明截面 $s_i$ 形变为 $s'_i\in\Gamma(\mathcal C',\allowbreak L')$. 因为 $\ker(A'\to\allowbreak A)=k$, 定义 $\mathcal C\hookrightarrow\mathcal C'$ 的理想层同构于 $\mathcal O_C$, 有短正合列
\[
0\to\mathcal O_C\to\mathcal O_{\mathcal C'}\to\mathcal O_{\mathcal C}\to0.
\]
张量 $L'$ 得到 $0\to\allowbreak L|_C\to\allowbreak L'\to\allowbreak L\to0$. 由于 $[C\subseteq\mathbb P^N]\in H'$, 有 $\mathrm H^1(C,\allowbreak L|_C)=0$. 因而 $\mathrm H^0(\mathcal C',\allowbreak L')\to\mathrm H^0(\mathcal C,\allowbreak L)$ 满射, 可将 $s_i$ 提升为 $s'_i$. 截面 $s_i$ 无基点, Nakayama 引理说明 $s'_i$ 也无基点. 由此得到限制为 $\mathcal C\hookrightarrow\mathbb P^N_A$ 的态射 $j':\mathcal C'\to\mathbb P^N_{A'}$. 映射 $j'$ 固有且拟有限, 因而有限. 为说明它是闭浸入, 只需证明 $\mathcal O_{\mathbb P^N_{A'}}\to\allowbreak j'_*\mathcal O_{\mathcal C'}$ 满射. 余核是限制到 $A$ 后为零的凝聚层, 由 Nakayama 引理余核为零. 另见~\cite[定理 B.1]{Hal13} 和~\cite[命题 3.3]{JHS11}. 也可用 Artin 公理建立此结果; 见定理~\bookref{C.7.7} 和~\cite[Tag 0D5A]{SP}.
\end{proof}
\begin{exercise}[容易]\label{src:6.4.8}
证明忘掉最后一个标记点的映射 $\mathcal M_{g,\allowbreak n+1}^{\mathrm{all}}\to\mathcal M_{g,\allowbreak n}^{\mathrm{all}}$ 是定义~\bookref{4.1.29} 意义下的万有族. (命题~\bookref{6.6.12} 将说明 $\overline{\mathcal M}_{g,\allowbreak n+1}$ 也是 $\overline{\mathcal M}_{g,\allowbreak n}$ 的万有族; 这更为惊人, 因为 $n$-点稳定曲线在忘掉一个标记点后可能变得不稳定.)
\end{exercise}
\begin{exercise}[技术性]\label{src:6.4.9}
证明 $\mathcal M_{g,\allowbreak n}^{\mathrm{all}}$ 的对角线拟紧且分离, 特别地 $\mathcal M_{g,\allowbreak n}^{\mathrm{all}}$ 拟分离.
\end{exercise}
\begin{exercise}[很好的练习]\label{src:6.4.10}
证明参数化曲线族 $C\to\allowbreak S$ 及 $C$ 上在 $S$ 上相对丰沛的线丛 $L$ 的叠 $\mathcal P$ 是代数叠, 且自然映射 $\mathcal P\to\mathcal M_g^{\mathrm{all}}$ 光滑且可表.
\end{exercise}
\begin{remark}\label{src:6.4.11}
更高维簇的模叠常常不是代数叠. 例如参数化抽象 K3 曲面的叠不是代数叠 (见例~\bookref{C.7.15}); 另一方面, 有参数化 K3 曲面的解析叠 (见~\cite[\S VII.12]{BHPV04}), 以及参数化极化 K3 曲面 (即配备本原丰沛线丛的 K3 曲面) 的代数叠.
\end{remark}
\begin{exercise}\label{src:6.4.12}
证明对所有 $g$, $\mathcal M_g^{\mathrm{all}}$ 连通.

\textit{提示:} 用 Hilbert 概形 $\Hilb^P(\mathbb P^N)$ 的连通性; 见 \S\bookref{2.5.2-subsection}.
\end{exercise}
\begin{remark}\label{src:6.4.13}
更有意思的是, 连通约化曲线的子叠连通~\cite{Boz25}. 另一方面, Mumford~\cite{Mum75b} 和 Pinkham~\cite[定理 1.11]{Pin74} 的例子说明 $\mathcal M_g^{\mathrm{all}}$ 有完全由奇异曲线组成的不可约分支.
\end{remark}
\subsection{稳定曲线模叠的代数性与有界性}\label{src:6.4.3-subsection}
考虑预叠的包含
\begin{equation}\label{src:6.4.14}
\mathcal M_{g,n}\subseteq\overline{\mathcal M}_{g,n}\subseteq\mathcal M_{g,n}^{\mathrm{ss}}\subseteq\mathcal M_{g,n}^{\mathrm{pre}}\subseteq\mathcal M_{g,n}^{\leq\mathrm{nodal}}\subseteq\mathcal M_{g,n}^{\mathrm{all}},
\end{equation}
其中 $\overline{\mathcal M}_{g,\allowbreak n}$ (分别地, $\mathcal M_{g,\allowbreak n}^{\mathrm{ss}}$, $\mathcal M_{g,\allowbreak n}^{\mathrm{pre}}$, $\mathcal M_{g,\allowbreak n}^{\leq\mathrm{nodal}}$) 是 $\mathcal M_{g,\allowbreak n}^{\mathrm{all}}$ 中由稳定 (分别地, 半稳定, 预稳定, 结点) 曲线的 $n$-点族 $(C\to\allowbreak S,\allowbreak \sigma_1,\allowbreak \ldots,\allowbreak \sigma_n)$ 组成的全子范畴.
\begin{corollary}\label{src:6.4.15}
\bookeqref{6.4.14} 中的包含都是开浸入, 每个预叠都是 $\Spec\mathbb Z$ 上局部有限型的代数叠.
\end{corollary}
\begin{proof}
定理~\bookref{6.4.6} 说明 $\mathcal M_{g,\allowbreak n}^{\mathrm{all}}$ 代数且局部有限型. 命题~\bookref{6.3.24} 已证明稳定 (分别地, 半稳定, 预稳定, 结点) 轨迹的开放性, 因而每个包含都是开浸入.
\end{proof}
至此已知 $\overline{\mathcal M}_{g,\allowbreak n}$ 的以下性质.
\begin{theorem}\label{src:6.4.16}
假设 $2g-2+n>0$, 则叠 $\overline{\mathcal M}_{g,\allowbreak n}$ 非空, 拟紧且为 Deligne--Mumford 叠, 在 $\Spec\mathbb Z$ 上光滑, 对每个域 $k$, $\overline{\mathcal M}_{g,\allowbreak n}\times_{\mathbb Z}k$ 的纯维数为 $3g-3+n$.
\end{theorem}
\begin{proof}
由推论~\bookref{6.4.15}, $\overline{\mathcal M}_{g,\allowbreak n}$ 是 $\Spec\mathbb Z$ 上有限型代数叠. 由于任意亏格都有光滑曲线 (习题~\bookref{6.1.7}), $\overline{\mathcal M}_{g,\allowbreak n}$ 非空. 为证明有界性 (即有限型, 或等价地拟紧性), 使用以下事实: 若 $(C,\allowbreak p_1,\allowbreak \ldots,\allowbreak p_n)$ 是域 $k$ 上的 $n$-点稳定曲线, 则 $L:=(\omega_{C/k}(p_1+\cdots+p_n))^{\otimes3}$ 甚丰沛 (习题~\bookref{6.3.18}). 设 $P(t)$ 是通过 $L$ 嵌入的 $C\hookrightarrow\mathbb P^N_k$ 的 Hilbert 多项式; 它与 $[C,\allowbreak p_i]\in\overline{\mathcal M}_{g,\allowbreak n}$ 无关. 考虑闭子概形
\[
H\subseteq\Hilb^P(\mathbb P^N_{\mathbb Z}/\mathbb Z)\times(\mathbb P^N)^n,
\]
参数化满足 $p_i\in C$ 的 $(C\hookrightarrow\mathbb P^N,\allowbreak p_1,\allowbreak \ldots,\allowbreak p_n)$. 因为 $\Hilb^P(\mathbb P^N_{\mathbb Z}/\mathbb Z)$ 是射影概形 (定理~\bookref{2.1.2}), 特别地拟紧, $H$ 也拟紧. 忘却态射 $H\to\mathcal M_{g,\allowbreak n}^{\mathrm{all}}$, $[C\hookrightarrow\mathbb P^N,\allowbreak p_1,\allowbreak \ldots,\allowbreak p_n]\mapsto[C,\allowbreak p_1,\allowbreak \ldots,\allowbreak p_n]$ 的像包含 $\overline{\mathcal M}_{g,\allowbreak n}$, 所以 $\overline{\mathcal M}_{g,\allowbreak n}$ 拟紧.

为证明其余断言, 应用命题~\bookref{6.3.26} 各项对稳定曲线的自同构, 形变和障碍的刻画. 这类似于命题~\bookref{4.7.7} 的证明, 它对光滑曲线的叠 $\mathcal M_g$ 断言相同性质. 设 $(C,\allowbreak p_i)$ 是亏格为 $g$ 的 $n$-点稳定曲线. 由于 $\operatorname{Ext}^0_{\mathcal O_C}(\Omega_C(\sum_i p_i),\allowbreak \mathcal O_C)=0$, $\underline{\operatorname{Aut}}(C,\allowbreak p_i)$ 的 Lie 代数平凡, 且它是有限约化群概形. 由 Deligne--Mumford 叠的刻画 (\bookref{4.6.4}), $\overline{\mathcal M}_{g,\allowbreak n}$ 是 Deligne--Mumford 的. 由于 $\operatorname{Ext}^2_{\mathcal O_C}(\Omega_C(\sum_i p_i),\allowbreak \mathcal O_C)=0$, 稳定曲线的形变没有障碍; 无穷小提升判据 (\bookref{4.7.1}) 说明 $\overline{\mathcal M}_{g,\allowbreak n}$ 在 $\Spec\mathbb Z$ 上光滑. 最后, $[C,\allowbreak p_i]\in\overline{\mathcal M}_{g,\allowbreak n}\times_{\mathbb Z}k$ 的 Zariski 切空间与 $\operatorname{Ext}^1_{\mathcal O_C}(\Omega_C(\sum_i p_i),\allowbreak \mathcal O_C)$ 双射, 后者维数为 $3g-3+n$, 故 $\overline{\mathcal M}_{g,\allowbreak n}\to\Spec\mathbb Z$ 的相对维数为 $3g-3+n$.
\end{proof}
\begin{exercise}\label{src:6.4.17}
将 $\overline{\mathcal M}_{g,\allowbreak n}$ 明确表为商叠 $[H'/\operatorname{PGL}_N]$, 其中 $H'$ 是拟射影概形, 证明它的代数性.

\textit{提示:} 仿照定理~\bookref{4.1.17} 的证明.
\end{exercise}
\begin{exercise}\label{src:6.4.18}
证明 $\overline{\mathcal M}_{g,\allowbreak n}$ 拟分离.

\textit{提示:} 用前一习题说明其对角线事实上仿射. 或者为证明 $T$ 上的两个稳定曲线族 $C_1,\allowbreak C_2$ 的 $\operatorname{Isom}_T(C_1,\allowbreak C_2)\to\allowbreak T$ 拟紧, 利用同构的图是 $C_1\times_T C_2$ 的闭子概形, 及相对 Hilbert 概形 $\Hilb^P(C_1\times_T C_2/T)$ 的有界性.
\end{exercise}
\begin{exercise}\label{src:6.4.19}
证明 $\mathcal M_{g,\allowbreak n}^{\mathrm{ss}}$, $\mathcal M_{g,\allowbreak n}^{\mathrm{pre}}$, $\mathcal M_{g,\allowbreak n}^{\leq\mathrm{nodal}}$ 和 $\mathcal M_{g,\allowbreak n}^{\mathrm{all}}$ 都不拟紧.

\textit{提示:} 利用任意长度的有理桥的存在性.
\end{exercise}
\begin{remark}\label{src:6.4.20}
$\mathcal M_g^{\mathrm{all}}$ 还有许多参数化特定曲线类的开子叠, 例如
\[
\begin{aligned}
\mathcal M_g^{\mathrm{DM}}&=\{\text{具有有限约化自同构群的曲线}\},\\
\mathcal M_g^{\mathrm{CM}}&=\{\text{Cohen--Macaulay 曲线}\},\\
\mathcal M_g^{\mathrm{geom\text{-}red}}&=\{\text{几何约化曲线}\},\\
\mathcal M_g^{\mathrm{Gor}}&=\{\text{Gorenstein 曲线}\},\\
\mathcal M_g^{\mathrm{lci}}&=\{\text{局部完全交曲线}\},\\
\mathcal M_g^{\mathrm{iso}}&=\{\text{具有孤立奇点的曲线}\};
\end{aligned}
\]
见~\cite[Tags 0DST--0E0K]{SP}. 开子叠 $\mathcal M_g^{\mathrm{DM}}$ 等同于 $\mathcal M_g^{\mathrm{all}}$ 的最大 Deligne--Mumford 开子叠. 对任意 $g$, $\mathcal M_g^{\mathrm{all}}\setminus\mathcal M_g^{\mathrm{DM}}$ 中都有自同构群正维且非约化的曲线. 虽然 $\mathcal M_{g,\allowbreak n}^{\mathrm{all}}$ 在 $\Spec\mathbb Z$ 上不光滑, 开子叠 $\mathcal M_g^{\mathrm{lci}}\cap\mathcal M_g^{\mathrm{iso}}$ 光滑~\cite[Tag 0DZX]{SP}. 开子叠 $\mathcal M_g\subseteq\mathcal M_g^{\mathrm{all}}$ 不稠密: Mumford 给出了不能形变为光滑曲线的约化不可约固有曲线的例子~\cite{Mum75b}. 然而命题~\bookref{6.7.26} 将证明 $\mathcal M_g$ 在 $\overline{\mathcal M}_g$ 中稠密; 它在 $\mathcal M_g^{\mathrm{lci}}\cap\mathcal M_g^{\mathrm{iso}}$ 中也稠密~\cite[Tag 0E86]{SP}.
\end{remark}
\begin{exercise}[不重要]\label{src:6.4.21}
证明 $\mathcal M_g^{\mathrm{all}}\to\Spec\mathbb Z$ 满足固有性的赋值判据的存在性部分, 但不普遍闭.

\textit{提示:} 考虑 $\mathcal M_g^{\mathrm{all}}\times\mathbb A^1_{\mathbb Z}$ 的闭子叠, 它由正整数 $n$ 所对应的 $([C_n],\allowbreak n)$ 的不交并组成, 其中 $C_n$ 是亏格 $g$ 曲线与具有 $n$ 个结点的 $\mathbb P^1$ 树的结点并. 该论证还应说明 $\mathcal M_g^{\mathrm{pre}}$ 在任意域上都不普遍闭.
\end{exercise}
